\title{Auditable by Design: Secure AI-Assisted Migration\\ to Post-Quantum Cryptography}

\author{Sergey~Mkrtchyan%
\thanks{Manuscript prepared October 2026; regulatory references reflect the policy context available on October 1, 2026.}%
\thanks{S.~Mkrtchyan is an independent researcher in Yerevan, Armenia (e-mail: b1acksparrow@hotmail.com; ORCID: 0009-0007-8442-1409).}%
}

\markboth{}{Mkrtchyan: Auditable by Design}

\maketitle

\begin{abstract}
An artificial intelligence (AI) recommendation for post-quantum migration neither authorizes a configuration change nor shows that it occurred, and explaining a decision based on a confidential rule discloses it. Auditable by Design is a transaction contract for cryptographic-configuration changes. A policy gate, re-evaluated before execution, checks each AI proposal against a signed observation; a one-time approval of the exact action permits a compare-and-set change, observed separately and recorded as post-quantum-signed, hash-chained, witnessed evidence. A streaming verifier enforces thirteen cross-statement obligations that reject correctly signed but inconsistent evidence from valid key holders, and an optional relation binds a confidential policy predicate to the same action. The single advisor model evaluated, Claude Sonnet~5 through Claude Code, followed \AdvHardViolFollowed\ planted policy-violating instructions under a hardened prompt and \AdvNaiveViolFollowed\ under a naive one; under complete proposer compromise, \AdvCompViolExecuted\ such proposals executed. Within-policy manipulation, which the gate permits, executed in \AdvHardInPolicyExecuted\ hardened, \AdvNaiveInPolicyExecuted\ naive, and \AdvCompInPolicyExecuted\ compromised trials. A RISC Zero succinct proof of the relation takes \ZkProveMed~s to generate and \ZkVerifyMs~ms to verify, with \ZkEndBits~bits of soundness under the vendor's toy-model conjecture and an unproven zero-knowledge property. Outside the contract, a four-client panel detects two edits of a Transport Layer Security~1.3 service, each silently downgrading one hybrid client class to classical key exchange; no single class detects both. On one host, a governed transaction takes a steady-state median of \TimDurTotal~ms (\TimDurSqlitePctR\% in durable commits), and a 100,000-transaction archive verifies in \ScalHundredKMedR~s within \ScalHundredKRss~MiB.
\end{abstract}

\begin{IEEEkeywords}
Artificial intelligence, audit trails, cryptographic agility, large language models, migration, post-quantum cryptography, security assurance, zero-knowledge proofs.
\end{IEEEkeywords}

\section{Introduction}
Post-quantum cryptography (PQC) is becoming an obligation in government and regulated environments. In the United States, Executive Order 14412 directs the Director of the Office of Management and Budget (OMB) to issue guidance requiring agencies to transition federal high-value assets and high-impact systems to PQC key establishment by December 31, 2030, and to PQC signatures by December 31, 2031. It also orders a proposed rule that would require covered contractors to comply with the applicable National Institute of Standards and Technology (NIST) Federal Information Processing Standards (FIPS) by December 31, 2030~\cite{eo14412}. OMB Memorandum M-26-15 implements the federal program~\cite{m2615}. The National Security Agency (NSA) states that Committee on National Security Systems Policy~15 (CNSSP~15) requires new commercial national security systems to support the Commercial National Security Algorithm Suite~2.0 (CNSA~2.0) by 2027, with legacy systems phased out by 2030~\cite{nsaarmor}. The European Commission targets the end of 2030 for protecting critical infrastructure under a coordinated roadmap for Member States~\cite{euroadmap,eunews}, and the United Kingdom's National Cyber Security Centre recommends completing the highest-priority activities by 2031 and migration by 2035~\cite{ncsc}. Organizations must turn applicable deadlines into executable changes and retain evidence of what changed, under whose authority, and with what outcome.

In FIPS 203 and 204, NIST standardized the Module-Lattice-Based Key-Encapsulation Mechanism (ML-KEM) and the Module-Lattice-Based Digital Signature Algorithm (ML-DSA)~\cite{fips203,fips204}. A migration also changes protocols, certificates, hardware security modules, dependencies, and recovery procedures~\cite{ir8547,cswp39}. Artificial intelligence (AI) can assist with discovery, dependency analysis, and configuration proposals, but its recommendations can rest on stale inventories, and large language model (LLM) agents given ordinary engineering prose have produced validated post-quantum downgrades~\cite{han}.

A recommendation establishes neither permission to change a system nor evidence that the change occurred, and such evidence must remain checkable against holders of valid signing keys. Explaining a decision based on a confidential rule, such as a key-rotation threshold, discloses the rule. Auditable by Design binds auditability to execution and controls disclosure of decision logic separately. Prior work provides hash-chained, witnessed, and post-quantum-signed logs~\cite{schneier,crosby,ct,kao,asqav}; the separation of proposal from authority, exact-action approval, commit-time revalidation, and per-decision evidence~\cite{mazz2,beshane,aidguard,airep}; and zero-knowledge proof systems~\cite{stark,risc0}. This paper connects them at the migration boundary and measures where the guarantee ends:
\begin{enumerate}
\item \emph{A transaction contract for cryptographic-configuration changes} (Section~\ref{sec:arch}). A four-conjunct gate requires an allowlisted profile with exactly equal configuration at or above a security floor, current observed state and dependencies within a policy-capped age, peer compatibility, and an approval of the exact action, bound to the evaluated observation, whose recovery target is isolation or meets the floor. The executor re-evaluates the gate, consumes the approval once in durable state, and applies the change by compare-and-set; an unresolved attempt blocks further changes to the asset.
\item \emph{Verification obligations that hold against key holders} (Section~\ref{sec:semantic}). Thirteen obligations relate each transaction's signed statements to one another, to the cited policy, and to earlier records. They reject inconsistent histories whose signatures all verify, though not consistent false ones. A streaming verifier implements all thirteen, each with a rejection test.
\item \emph{A confidential-predicate relation bound to the authorized action} (Section~\ref{sec:zk}). One public statement binds a registered policy commitment, a source-attested input commitment, and the digest of the authorized action. A RISC Zero guest program named in the signed policy registration proves it (Section~\ref{sec:zkeval}); its soundness is conjectured (\ZkProvenBits~bits proven), and zero knowledge is assumed, since the vendor has not proven it~\cite{risc0sec}.
\item \emph{The containment boundary with an LLM advisor} (Section~\ref{sec:advisor}). Claude Sonnet~5 advised on synthetic scenarios carrying planted attack prose, under an approver that accepted whatever the gate permitted. Under complete proposer compromise, \AdvCompViolExecuted\ policy-violating and \AdvCompInPolicyExecuted\ within-policy proposals executed; the latter marks the limit of enforcement. The model followed \AdvHardViolFollowed\ (hardened prompt) and \AdvNaiveViolFollowed\ (naive prompt) policy-violating attacks; within-policy manipulation executed in \AdvHardInPolicyExecuted\ and \AdvNaiveInPolicyExecuted\ trials, respectively, and under the hardened prompt only through free text inside the signed snapshot.
\item \emph{Wire-level observation of a deployed service} (Section~\ref{sec:tls}). Outside the contract, a panel of four reference clients reads the key-exchange group that a Transport Layer Security (TLS)~1.3 service negotiates with each. Two group-list edits each validate and keep serving but silently move a different hybrid client class to a classical group; only the panel detects both, consistent with Han's single-peer argument~\cite{han}.
\end{enumerate}
